\documentclass[12pt,reqno]{amsart}
\usepackage{amsmath}
\usepackage{amsfonts}
\usepackage{amssymb}
\usepackage{mathtools}
\usepackage{graphicx}
\usepackage{color}
\usepackage{hyperref}
\usepackage{geometry}
\usepackage{tikz-cd}
\parindent 15pt
\parskip 4pt

\usepackage{enumerate}

\textwidth=6.5in
\textheight=8.7in
\topmargin=0in
\oddsidemargin=0in
\evensidemargin=0in

\usepackage{xcolor}
\newtheorem{theorem}{Theorem}[section]
\newtheorem*{theorem*}{Theorem}
\newtheorem{lemma}[theorem]{Lemma}
\newtheorem{corollary}[theorem]{Corollary}
\newtheorem{proposition}[theorem]{Proposition}
\newtheorem{conjecture}[theorem]{Conjecture}

\theoremstyle{definition}
\newtheorem{definition}[theorem]{Definition}
\newtheorem{example}[theorem]{Example}
\newtheorem{exercise}[theorem]{Exercise}
\newtheorem{problem}[theorem]{Problem}
\newtheorem{remark}{Remark}[section]
\newtheorem{question}{Question}[section]
\renewcommand*{\thequestion}{\Alph{question}}


\DeclareMathOperator{\grad}{grad}
\DeclareMathOperator{\dimension}{dim}
\DeclareMathOperator{\trace}{tr} 
\DeclareMathOperator{\di}{div}
\DeclareMathOperator{\ricci}{Ricci}
\DeclareMathOperator{\dist}{dist}
\DeclareMathOperator{\hess}{Hess}

\def\zz{{\mathbb Z}}
\def\rr{{\mathbb R}}
\def\cc{{\mathbb C}}
\def\qq{{\mathbb Q}}
\def\oo{{\infty}}

\def\st{{S_{\mathcal{H},\widetilde{\mathcal{H}}}(f)}}
\def\dt{{df_{\mathcal{H},\widetilde{\mathcal{H}}}}}
\def\dtu{{du_{\mathcal{H},\widetilde{\mathcal{H}}}}}
\def\foli{{(M^{m+k},\mathcal{F}^k,g)}}
\def\folin{{(N^{n+l}, \widetilde{\mathcal{F}}^l , \widetilde{g})}}
\def\HH{{\mathcal{H},\widetilde{\mathcal{H}}}}

\newcommand{\la}{\langle}
\newcommand{\ra}{\rangle}


	
\allowdisplaybreaks
	
	
\begin{document}
    \begin{abstract}
        In this paper, we consider critical points of the
        horizontal energy $E_{\HH}(f)$ for a smooth map $f$ between two Riemannian
        foliations. These critical points are referred to as horizontally harmonic
        maps. In particular, if the maps are foliated, they become transversally
        harmonic maps. By utilizing the stress-energy tensor, we establish some
        monotonicity formulas for horizontally harmonic maps from Euclidean spaces,
        the quotients $K_{m}$ of Heisenberg groups and also for transversally
        harmonic maps from Riemannian foliations with appropriate curvature pinching
        conditions. Finally, we give Jin-type theorems for either horizontally
        harmonic maps or transversally harmonic maps under some asymptotic
        conditions at infinity.
\end{abstract} %%%%%%%%%


    
    \title[Horizontal maps between Riemannian foliations]{On critical maps of the horizontal energy functional between Riemannian foliations}

    \author{Tian Chong}
    \date{\today}
    \address{Department of Mathematics,
Shanghai Polytechnic University, Shanghai 201209, PR China}
    \email{chongtian@sspu.edu.cn}
    
    \author{Yuxin Dong}
    \address{School of Mathematical Science  and  Laboratory of Mathematics for Nonlinear Science, Fudan University, Shanghai 200433, PR China}
    \email{yxdong@fudan.edu.cn}
    
    \author{Xin Huang }
    \address{School of Mathematics and Statistics, Nanjing University of Information Science and Technology,
Nanjing 210044,
PR China}
    \email{17110180003@fudan.edu.cn}
    
    \author{Hui Liu}
    \address{School of Mathematical Sciences, Fudan University, Shanghai 200433, PR China}
    \email{21110180015@m.fudan.edu.cn and liuh45@univie.ac.at}
    \maketitle
    
    \let\thefootnote\relax
    \footnotetext{ \textbf{Mathematics Subject Classification} 58E20, 53C12, 53C43 } %%%%%%%%%%		

  
\section{Introduction}

    Harmonic maps between Riemannian manifolds are an important object of study in
	geometric analysis. Over the decades, several generalized harmonic maps have
	appeared in the literature, playing an important role in geometry and
	topology. These generalizations include some generalized harmonic maps from
	Riemannian foliations, CR manifolds and contact manifolds (see \cite{konderakREMARKSTRANSVERSALLYHARMONIC2008, konderakTransversallyHarmonicMaps2003}, \cite{chenBochnertypeFormulasTransversally2012a}, 
	\cite{dragomirHarmonicMapsFoliated2013a},
	\cite{barlettaTransversallyHolomorphicMaps1998a}, \cite{barlettaPseudoharmonicMapsNondegenerate2001},
	\cite{petitMokSiuYeungTypeFormulas2009},
	\cite{dragomirLeviHarmonicMaps2014},
	\cite{renPseudoharmonicMapsClosed2018}, \cite{chongHARMONICPSEUDOHARMONICMAPS2019},
	etc.).
	
	Riemannian foliations are a natural generalization of Riemannian manifolds.
	Let $(M,g,\mathcal{F})\ $ and $(\widetilde{M},\widetilde{g},\widetilde{%
	\mathcal{F}})$ be two Riemannian foliations and $f:M\rightarrow \widetilde{M}
	$ be a smooth map. We say that $f$ is foliated if it maps leaves of $%
	\mathcal{F}$ to leaves of $\widetilde{\mathcal{F}}$. Then a foliated map $f$
	induces a map $\widehat{f}:M/\mathcal{F\rightarrow }\widetilde{M}/\widetilde{%
	\mathcal{F}}$ between the leaf spaces. Since the leaf spaces can be very
	complicated topological spaces, it is difficult to investigate $\widehat{f}$
	directly. One has to study related problems about the map $f$ on the upstairs
	spaces. In order to generalize the theory of harmonic maps to Riemannian
	foliations, Konderak and Wolak (cf. \cite{konderakTransversallyHarmonicMaps2003,konderakREMARKSTRANSVERSALLYHARMONIC2008}) constructed a global
	transversal tension field for a foliated map $f$ in terms of the transversal
	second fundamental form, and then they defined transversally harmonic maps
	as those smooth maps for which this tension field vanishes. However, S.
	Dragomir and A. Tommasoli (\cite{dragomirHarmonicMapsFoliated2013a}) pointed out that such maps do not, in
	general, extremize the natural transverse energy functional that will be
	described below.
	
	Let $\mathcal{V}(\mathcal{F})$ and $\mathcal{V}(\widetilde{\mathcal{F}})$ be
	the distributions that are determined by the foliations $\mathcal{F}$ and $%
	\widetilde{\mathcal{F}}$ respectively. Then we have the quotient bundles $Q(%
	\mathcal{F}):=T(M)/\mathcal{V}(\mathcal{F})$ and $Q(\widetilde{\mathcal{F}}%
	):=T(\widetilde{M})/\mathcal{V}(\widetilde{\mathcal{F}})$, which are called
	the transverse bundles of the foliations. These transverse bundles get
	natural fiberwise metrics from $g$ and $\widetilde{g}$ respectively. The
	differential map $df$ of $f$ induces the transverse differential $d_{T}f:Q(%
	\mathcal{F})\rightarrow f^{-1}(Q(\widetilde{\mathcal{F}}))$ in a natural
	way. Consequently we may define the norm $\Vert d_{T}f\Vert $ and the
	following transverse energy functional $E_{T}(f)$ (cf. \cite{dragomirHarmonicMapsFoliated2013a} for details): 
	\begin{equation}\label{tag1}
	E_{T}(f)=\int_{M}e_{T}(f)dV_{g}.  
	\end{equation}%
	It turns out that the transversally harmonic maps defined by Konderak and
	Wolak are not the critical points of the functional $E_{T}(f)$, but rather
	the critical points of the following energy functional%
	\begin{equation}\label{tag2}
	E_{T}^{\ast }(f)=\int_{M}\frac{e_{T}(f)}{vol_{L}}dV_{g} , \tag{2}
	\end{equation}%
	where $vol_{L}$ denotes the volume of each fiber. In \cite{dragomirHarmonicMapsFoliated2013a}, S. Dragomir and
	A. Tommasoli proposed a different definition of generalized harmonic maps,
	which are exactly the extremals of $E_{T}(f)$ through any variation of $f$
	by foliated maps. They call such maps $(\mathcal{F},\widetilde{\mathcal{F}})$%
	-harmonic maps. It is easy to verify that the above two definitions of
	generalized harmonic maps coincide when the foliation $\mathcal{F}$ is
	harmonic (cf. \cite{chenBochnertypeFormulasTransversally2012a}, \cite{dragomirHarmonicMapsFoliated2013a}).
	
	In this paper, we introduce the horizontal energy functional $E_{\mathcal{H},%
	\widetilde{\mathcal{H}}}(f)$ similar to \eqref{tag1} for any map $f:(M^{m+k},\mathcal{%
	F},g)\rightarrow (N^{n+l},\widetilde{\mathcal{F}},\widetilde{g})$ (not
	necessarily foliated) between the Riemannian foliations, and then define
	horizontally harmonic maps as the critical points of $E_{\mathcal{H},%
	\widetilde{\mathcal{H}}}(f)$ through any variation (see \S \ref{se:criticalmaps} for details).
	The notion of horizontally harmonic maps slightly generalizes the notion of $%
	(\mathcal{F},\widetilde{\mathcal{F}})$-harmonic maps in \cite{dragomirHarmonicMapsFoliated2013a}. Notice that
	removing the foliated condition about the map in our definition allows the
	notion to involve the subelliptic harmonic maps when the domain is a
	Riemannian foliation and the target is a Riemannian manifold with the
	trivial foliation by points. The main purpose of the present paper is to
	study the properties of the horizontally harmonic maps, especially the
	Liouville type theorems for them.
	
	As known, the stress-energy tensors are a useful tool for deriving
	monotonicity formulas of harmonic maps or their generalizations (cf. \cite{price1983amono},
	\cite{xin1986diff}, \cite{dongVanishingTheoremsVector2011a}, \cite{dlyLiou2016}, and the references therein). Since the stress-energy tensors
	are naturally linked to the conservation laws of related energy functionals,
	the monotonicity formulas of the energies follow from Stokes' theorem,
	coarea formula and comparison theorems in Riemannian geometry. These
	monotonicity formulas can be used to establish Liouville type theorems in
	the following two ways. One way is to derive Liouville type theorems
	directly from the monotonicity formulas by assuming the energy growth
	conditions (cf. \cite{sealeysome1982}, \cite{hu1984anonexist}, \cite{ll2004theenergy}, \cite{dongVanishingTheoremsVector2011a}, \cite{rs2000energy}, and the references therein).
	Another way is to establish Liouville type theorems by assuming suitable
	asymptotic conditions of the maps at infinity (cf. \cite{jin_liouville_nodate}, \cite{dlyLiou2016}, \cite{rs2000energy}).
	
	For our purpose, we introduce the stress-energy tensor $S_{\mathcal{H},%
	\widetilde{\mathcal{H}}}(f)$ associated with the energy functional $E_{%
	\mathcal{H},\widetilde{\mathcal{H}}}(f)$ and derive the corresponding
	conservation law formula for $(\di S_{\mathcal{H},\widetilde{\mathcal{H
    	}}}(f))(X)$ with respect to any vector field $X\in \mathfrak{X}(M)$. It
	turns out that a horizontally harmonic map $f:(M^{m+k},\mathcal{F}
	,g)\rightarrow (N^{n+l},\widetilde{\mathcal{F}},\widetilde{g})$ does not
	satisfy the conservation law in general, but if $f$ is a transversally
	harmonic map, then it satisfies the conservation law. Even though in the
	general case $f$ does not satisfy the conservation law, we find that the
	conservation law formula becomes more manageable if  the domain
	Riemannian foliation is simple. To illustrate this, we consider the
	following two cases for horizontally harmonic maps: (i) $(M^{m+k},\mathcal{F%
	},g)=(%
	\mathbb{R}
	^{m+k},%
	\mathbb{R}
	^{k},g)$ the Euclidean space foliated by Euclidean subspaces, where $g$ is
	some mixed conformally flat metric on $%
	\mathbb{R}
	^{m+k}$ (see Example \ref{exam4}); (ii) $(M^{m+k},\mathcal{F},g)=(K_{m},S^{1},g)$
	the quotient space of the Heisenberg group $H_{m}$ with the canonical metric 
	$g$ (see Example \ref{ex:sasa}). The horizontal distribution of the former is
	integrable, while the horizontal distribution of the latter is
	non-integrable. They can be viewed respectively as the simplest models in
	the two cases. Since the conservation law formula becomes simpler in
	horizontal directions for these cases, it is more appropriate to apply the
	formula on cylindrical regions defined by the distance function from a fixed
	leaf. Consequently, we can establish some monotonicity formulas of the
	horizontal energy on cylindrical regions for horizontally harmonic maps from
	either $\rr^{m+k}$ or $K_{m}$. As mentioned above, a transversally harmonic
	map  $f:(M^{m+k},\mathcal{F},g)\rightarrow (N^{n+l},\widetilde{\mathcal{F}},%
	\widetilde{g})$ between two Riemannian foliations satisfies the conservation
	law, that is, $(\di S_{\mathcal{H},\widetilde{\mathcal{H}}}(f))(X)=0$
	in any direction $X$. This enables us to establish some monotonicity
	formulas of the horizontal energy on geodesic balls for transversally
	harmonic maps under suitable curvature pinching conditions of $(M,g)$.
	Obviously, under the assumption of appropriate energy growth conditions, the
	above monotonicity formulas for horizontally harmonic maps or transversally
	harmonic maps give us directly Liouville type theorems for these maps.
	
	In \cite{jin_liouville_nodate}, Jin established several interesting Liouville type theorems for
	harmonic maps from Euclidean spaces endowed with conformally flat metrics,
	under an asymptotic condition of these maps at infinity. A special case of
	his results is that if $u:(\rr^{m},g_{can})\rightarrow (N^{n},h)$ is a
	harmonic map, and $u(x)$ converges to a fixed point $p_{0}\in N$ as $|
	x| \rightarrow \infty $, then $u$ is a constant map. His method is based
	on the following two steps for a non-constant harmonic map: first, deriving
	a lower bound on the energy growth rate, and second, deriving an upper bound
	on the energy growth rate under an asymptotic condition. If the two growth
	rates are incompatible, this shows that $u$ can only be constant. Inspired
	by Jin's method, we next investigate Jin-type Liouville theorems for a
	horizontally harmonic map or transversally harmonic map $f:(M^{m+k},\mathcal{%
	F},g)\rightarrow (N^{n+l},\widetilde{\mathcal{F}},\widetilde{g})$ by
	assuming that the map approaches a fixed point $q$ at infinity. We will
	establish Jin-type theorems for the following two cases: (i) $f:(
	\mathbb{R}
	^{m+k},
	\mathbb{R}
	^{k},g)\rightarrow (N^{n+l},\widetilde{\mathcal{F}},\widetilde{g})$ is a
	horizontally harmonic map, where $g$ is a mixed conformally flat metric on $
	\mathbb{R}
	^{m+k}$; (ii) $f:(M^{m+k},\mathcal{F},g)\rightarrow (N^{n+l},\widetilde{%
	\mathcal{F}},\widetilde{g})$ is a transversally harmonic map, where $g$ is a
	complete Riemannian metric with radial curvature $K_{r}$ satisfying $-\alpha
	^{2}\leq K_{r}\leq -\beta ^{2}$, $\alpha ,\beta >0$ and $(m+k-1)\beta
	-2\alpha \geq 0$. Note that the monotonicity formulas we have already
	established for these two cases in \S 5 actually give lower bounds on the
	energy growth rate. This completes the first step in Jin's method. For the
	second step, we may assume that the image of the map $f$ near infinity is
	contained in a foliated coordinate chart $(U,\widetilde{\mathcal{F}}|_{U})$,
	which induces a natural Riemannian submersion $\pi _{U}:U\rightarrow 
	\mathcal{B}_{U}$. Clearly, estimating the horizontal energy of $f$ near
	infinity is equivalent to estimating the horizontal energy of $\pi _{U}\circ
	f$. Then we can apply Jin's method to give the upper bound on the energy
	growth rate. Therefore, if the two growth rates are incompatible, we can
	establish a Jin-type Liouville theorem (see \S \ref{se:Jintypethm} for details).
	

This paper is structured in the following manner. 
In \S \ref{se:Riemfoli}, we collect some basic notions, propositions and some examples in Riemannian foliations. 
In \S \ref{se:criticalmaps}, we introduce the horizontal energy $E_{\mathcal{H}, \widetilde{\mathcal{H}}}(f)$ for a smooth map $f$, and then deduce the first variation formula, and we also explore the geometric meaning of horizontally constant maps.
In \S \ref{se:streeenergy}, we derive the divergence formula for the stress-energy tensor $S_{\HH}(f)$ and introduce the associated concepts of conservation laws.   
In \S \ref{se:monoformula}, 
we establish monotonicity formulas for horizontally harmonic maps from Euclidean spaces and the quotient space $K_m$ of the Heisenberg group $H_m$, as well as for transversally harmonic maps from Riemannian foliations under certain curvature pinching conditions.
 Finally, in \S \ref{se:Jintypethm}, we present Jin-type theorems for horizontally harmonic maps from mixed conformally flat Euclidean spaces and transversally harmonic maps from  Riemannian foliations with a pole.




~
\section*{Acknowledgments}
T was supported by National Natural Science Foundation of China (NSFC) No. xxx, 
Y was supported by NSFC Grant No. xxxx,
X was supported by the xxx  
and H was supported by the China Scholarship Council (No. xxxx).
The authors would like to thank Prof. V. B. for his valuable comments.

\section*{Declarations}

\textbf{Conflict of interest} The authors have no relevant financial or non-financial interests to disclose. The authors
have no Conflict of interest to declare that are relevant to the content of this article.

\textbf{Data Availability Statement} The authors did not collect
or use any data for this work.

\bibliographystyle{alpha}
\bibliography{ref.bib}   

\end{document}
